\documentclass[../differential_geometry.tex]{subfiles}
\begin{document}
\section{Aula 15 - Bônus}
\subsection{Motivações}
\begin{itemize}
	\item Operadores Auto-adjuntos.
\end{itemize}
\subsection{Operadores Auto-adjuntos}
\begin{def*}
	Um \textbf{operador autoadjunto} é um operador \(A:X\rightarrow X\) com
	\[
		\langle Ax, y \rangle = \langle x, Ay \rangle, \quad x, y\in T_{p}S.\; \square
	\]
\end{def*}

Se \(\underline{x}:U\subseteq \mathbb{R}^{2}\rightarrow S\subseteq \mathbb{R}^{3}\) é uma parametrização local de S, tal que
\[
	N(\underline{x}(u, v))=\frac{\underline{x}_u(u, v)\times \underline{x}_v(u, v)}{\Vert \underline{x}_u(u, v)\times \underline{x}_v(u, v) \Vert}.
\]
Vimos que a derivada da aplicação de Gauss \(\mathrm{d}N_{p}:T_pS\rightarrow T_pS\) é um operador autoadjunto, tal que ele satisfaz
\[
	\langle \mathrm{d}N_{p}(w_1), w_2 \rangle = \langle w_1, \mathrm{d}N_p(w_2) \rangle, \quad \forall w_1, w_2\in T_pS.
\]
Seja \(\ell:V\rightarrow V\) um operador e \(\{v_1, v_2\}\) uma base de V; escreva
\begin{align*}
	 & \ell(v_1) = a_{11}v_1 + a_{21}v_2  \\
	 & \ell(v_1) = a_{12}v_1 + a_{22}v_2.
\end{align*}
Então, a matriz de \(\underline{\ell} \) na base \(\{v_1, v_2\}\); assim, se \(w = \alpha v_1 + \beta v_2\in V\)  e \(\underline{\ell }(w) = \gamma v_1 + \eta v_2\), temos
\[
	\begin{bmatrix}
		\gamma \\
		\eta
	\end{bmatrix} = A \begin{bmatrix}
		\alpha \\
		\beta
	\end{bmatrix}
\]
O determinante não depende da escolha da base de V, e vale \(\det{(A)} = a_{11}a_{22} - a_{12}a_{21}\), pois se \(\tilde{A}\) é a matriz de \(\underline{\ell }\) em relação a uma outra base de V, temos:
\[
	\tilde{A} = PAP^{-1}
\]
para alguma matriz invertível P, ou seja, \(\det{\tilde{A}} = \det{A}\). Por conta disso, escrevemos simplesmente \(\det{(\underline{\ell })=a_{11}a_{22}-a_{12}a_{21}}\)

Por essa mesma lógica, o traço \(a_{11} + a_{22}\) de A também não depende da escolha da base de V, e escrevemos
\[
	\mathrm{tr}(\underline{\ell}) = a_{11} + a_{22}.
\]

\begin{theorem*}
	Seja \(\underline{\ell }:V\rightarrow V\) um operador autoadjunto; existe uma base ortonormal \(\{e_1, e_2\}\) de V de autovetores de \(\underline{\ell }\):
	\begin{align*}
		 & \underline{\ell }(e_1) = \lambda_1e_1  \\
		 & \underline{\ell }(e_2) = \lambda_2e_2,
	\end{align*}
	onde \(\lambda_1\) e \(\lambda_{2}\) são os autovetores de \(\underline{\ell }\). Além disso, a matriz de \(\underline{\ell }\) na base \(\{e_1, e_2\}\) é
	\[
		\mathrm{Matriz}(\underline{\ell }) = \begin{bmatrix}
			\lambda_1 & 0      \\
			0         & \ell_2
		\end{bmatrix},
	\]
	e valem as relações
	\begin{align*}
		 & \det{(\underline{\ell })}=\lambda_1\lambda_2            \\
		 & \mathrm{tr}(\underline{\ell }) = \lambda_1 + \lambda_2.
	\end{align*}
\end{theorem*}
Vamos calcular \(\det{(\underline{\ell })},\; \mathrm{tr}(\underline{\ell })\) em relação a uma base qualquer de V, denotada por \(\{v_1, v_2\}\). Além disso, escreveremos os coeficientes da primeira forma fundamental como
\[
	E = \langle v_1, v_1 \rangle,\quad F = \langle v_1, v_2 \rangle \quad\&\quad G = \langle v_2, v_2 \rangle;
\]
e os da segunda forma fundamental como
\[
	\ell = \langle \underline{\ell }(v_1), v_1 \rangle,\quad m = \langle \ell (v_1), v_2 \rangle \quad\&\quad n = \langle \underline{\ell} (v_2), v_2 \rangle.
\]
Com isso, seja A a matriz de \(\underline{\ell }\) em relação a \(\{v_1, v_2\}\) com coeficientes
\[
	A = \begin{bmatrix}
		a_{11} & a_{12} \\
		a_{21} & a_{22}
	\end{bmatrix};
\]
então, segue que:
\begin{align*}
	 & \tikz[baseline=-0.5ex]\draw[black, fill=black, radius=1.5pt](0, 0)circle;\; \ell = \langle \underline{\ell }(v_1), v_1 \rangle = \langle a_{11}v_1 + a_{21}v_2, v_1 \rangle = a_{11}\langle v_1, v_1 \rangle + a_{21}\langle v_{2}, v_1 \rangle = a_{11}E + a_{21}F \\
	 & \tikz[baseline=-0.5ex]\draw[black, fill=black, radius=1.5pt](0, 0)circle;\; m = \langle \underline{\ell }(v_1), v_2 \rangle = \langle v_1, \underline{\ell }(v_2) \rangle = a_{11}F + a_{21}G = a_{12}E + a_{22}F,                                                  \\
	 & \tikz[baseline=-0.5ex]\draw[black, fill=black, radius=1.5pt](0, 0)circle;\; n = \langle \underline{\ell }(v_2), v_2 \rangle = a_{12}F + a_{22}G,
\end{align*}
ou seja,
\[
	\begin{bmatrix}
		\ell & m \\
		m    & n
	\end{bmatrix} = \begin{bmatrix}
		E & F \\
		F & G
	\end{bmatrix} \begin{bmatrix}
		a_{11} & a_{12} \\
		a_{21} & a_{22}
	\end{bmatrix},
\]
com determinante
\[
	\det{\underline{\ell }} = \frac{\ell n - m^{2}}{EG - F^{2}}.
\]
Temos, também,
\[
	\begin{bmatrix}
		a_{11} & a_{12} \\
		a_{21} & a_{22}
	\end{bmatrix} = \begin{bmatrix}
		E & F \\
		F & G
	\end{bmatrix}^{-1}\begin{bmatrix}
		\ell & m \\
		m    & n
	\end{bmatrix} = \frac{1}{EG-F^{2}}\begin{bmatrix}
		G  & -F \\
		-F & E
	\end{bmatrix}\begin{bmatrix}
		\ell & m \\
		m    & n
	\end{bmatrix};
\]
logo,
\[
	\mathrm{tr}(\underline{\ell }) = \frac{\ell G - 2mF + nE}{EG-F^{2}}.
\]

E quanto aos autovalores e autovetores de \(\underline{\ell }\)?
O escalar \(\lambda \) é um autovalor de \(\underline{\ell }\) se, e somente se,
\[
	\det \begin{pmatrix}
		a_{11} - \lambda & a_{12}           \\
		a_{21}           & a_{22} - \lambda
	\end{pmatrix} = 0 \Longleftrightarrow \lambda^{2}-\lambda \mathrm{tr}(\underline{\ell })+\det{(\underline{\ell })}=0.
\]
Portanto, \(\lambda_1,\; \lambda_2\) são soluções da equação
\[
	\lambda^{2} - \frac{\ell G - 2mF + nE}{EG-F^{2}}\lambda  + \frac{\ell n - m^{2}}{EG-F^{2}} = 0.
\]
Um vetor \(\alpha_1v_1+\alpha_2v_2\), por sua vez, será um autovetor de \(\underline{\ell }\) associados ao autovalor \(\lambda \) se, e somente se,
\begin{align*}
	                    & \underline{\ell }(\alpha_{1}v_1 + \alpha_2v_2) = \lambda (\alpha_1v_1+\alpha_2v_2)                     \\
	\Longleftrightarrow & \alpha_1(a_{11}v_1 + a_{21}v_2) + \alpha_2(a_{12}v_1 + a_{22}v_2) = \lambda(\alpha_1v_1 + \alpha_2v_2) \\
	\Longleftrightarrow & \left\{\begin{array}{ll}
		                             \alpha_1a_{11} + \alpha_2a_{12} = \lambda\alpha_1 \\
		                             \alpha_1a_{21} + \alpha_2a_{22} = \lambda\alpha_2
	                             \end{array}\right..
\end{align*}
Daí, eliminando \(\lambda\), obtemos
\[
	a_{12}\alpha_{1}^{2} + (a_{11}-a_{22})\alpha_1\alpha_2 - a_{21}\alpha_{1}^{2} = (Gm-Fn)\alpha_{2}^{2}+(G\ell -En)\alpha_{1}\alpha_{2}+(F\ell -Em)\alpha_{1}^{2}=0,
\]
que pode ser escrito como
\[
	\det \begin{bmatrix}
		\alpha_{2}^{2} & -\alpha_1\alpha_2 & \alpha_{1}^{2} \\
		E              & F                 & G              \\
		\ell           & m                 & n
	\end{bmatrix} = 0.
\]

Agora, associamos a \(\underline{\ell }:V\rightarrow V\) a forma quadrática
\begin{align*}
	Q: & V\rightarrow \mathbb{R}                                             \\
	   & w\longmapsto Q(w)\coloneqq \langle \underline{\ell }(w), w \rangle;
\end{align*}
com isto, seja \(\{e_1, e_2\}\) uma base ortonormal de autovetores de \(\underline{\ell }\) e w um vetor unitário em V, caracterizado por \(\langle w, w \rangle = 1\). Nestas condições, se \(\theta = \mathrm{ang}(w, e_1)\), obtemos
\[
	w = \cos^{}{\theta } e_1 + \sin^{}{\theta }e_2
\]
e
\begin{align*}
	Q(w) = \langle \underline{\ell }(w), w \rangle & =\langle \underline{\ell }(\cos^{}{\theta }e_1 + \sin^{}{\theta }e_2), \cos^{}{\theta }e_1 + \sin^{}{\theta }e_2 \rangle \\
	                                               & =\langle \cos^{}{\theta }\ell (e_1)+\sin^{}{\theta }\ell (e_2), \cos^{}{\theta }e_1 + \sin^{}{\theta }e_2 \rangle        \\
	                                               & =\lambda_1 \cos^{2}{\theta } + \lambda_2\sin^{2}{\theta },
\end{align*}
permitindo concluirmos que a forma quadrática associada a \(\underline{\ell }\) tem forma
\[
	\boxed{Q(w) = \lambda_1 \cos^{2}{\theta } + \lambda_2 \sin^{2}{\theta }}
\]
Derivando em relação a \(\theta \),
\[
	\frac{\mathrm{d}Q}{\mathrm{d}\theta }=2 (\lambda_2 - \lambda_1)\sin^{}{\theta }\cos^{}{\theta }.
\]
Portanto, se \(\lambda_1\neq \lambda_2\), Q atinge os valores máximos/mínimos quando \(w = \pm e_1\) ou \(w=\pm e_2\).

\end{document}

